%\documentclass[12pt,oneside]{amsart}
\documentclass{scrartcl}

%\documentclass{article}

\usepackage{amsthm,amsmath,amssymb}
\usepackage[numbers]{natbib}

%\usepackage[colorlinks]{hyperref}
\usepackage{hyperref}
\hypersetup{
    colorlinks,%
    citecolor=black,%
    filecolor=black,%
    linkcolor=black,%
    urlcolor=black
}
\usepackage{hypernat}
\usepackage[top=2.5cm, bottom=2.5cm, left=2.8cm,
right=2.8cm]{geometry}
\usepackage{graphicx}


%\setlength{\parskip}{0pt}
%\setlength{\parsep}{0pt}
%\setlength{\headsep}{0pt}
%\setlength{\topskip}{0pt}
%\setlength{\topmargin}{0pt}
%\setlength{\topsep}{0pt}
%\setlength{\partopsep}{0pt}

\linespread{1}

%\usepackage{marvosym}

%\textwidth = 6.5 in \textheight = 9.2 in \oddsidemargin = 0.0 in
%\evensidemargin = 0.0 in \topmargin = -0.0 in \headheight = 0.0 in
%\headsep = 0.0in \parskip = 0.0in \parindent = 0.0in
%\def\baselinestretch{0.871}


\newtheorem{theorem}{Theorem}[section]
\newtheorem{proposition}[theorem]{Proposition}
\newtheorem{problem}[theorem]{Problem}
\newtheorem{fact}[theorem]{Fact}
\newtheorem{lemma}[theorem]{Lemma}
\newtheorem{defn}[theorem]{Definition}
\newtheorem{corollary}[theorem]{Corollary}
\newtheorem{remark}{Remark}[section]
\newtheorem{example}[theorem]{Example}

\newcommand{\E}{{\mathbb E}}
\newcommand{\se}{{\mathcal{E}}}
\newcommand{\R}{{\mathbb R}}
\renewcommand{\P}{{\mathbb P}}
\newcommand{\ac}{{\mathcal{A}}}
\newcommand{\A}{{\mathcal{A}}}
\newcommand{\B}{{\mathcal{B}}}
\newcommand{\C}{{\mathcal{C}}}
\newcommand{\M}{{\mathcal{M}}}
\newcommand{\Mf}{{\mathfrak{M}}}
\newcommand{\T}{{\mathcal{T}}}
\newcommand{\Z}{{\mathcal{Z}}}
\newcommand{\K}{{\mathcal{K}}}
\newcommand{\lc}{{\mathcal{L}}}
\newcommand{\Ps}{{\mathcal{P}}}
\newcommand{\G}{{\mathcal{G}}}
\newcommand{\pa}{{\mathring{p}}}
\newcommand{\F}{{\cal F}}
\newcommand{\samp}{{\mathcal{X}}}
\newcommand{\X}{{\mathcal{X}}}
\newcommand{\hi}{{\hat{\Phi}}}
\newcommand{\data}{{\text{data}}}
\newcommand{\relu}{{\text{ReLU}}}

\newcommand{\ridge}{{\hat{\beta}_{\text{ridge}}(\lambda)}}
\newcommand{\lasso}{{\hat{\beta}_{\text{lasso}}(\lambda)}}
\newcommand{\ridgemu}{{\hat{\mu}_t^{\text{ridge}}(\lambda)}}
\newcommand{\lassomu}{{\hat{\mu}_t^{\text{lasso}}(\lambda)}}


\newcounter{rcnt}[section]
\renewcommand{\thercnt}{(\roman{rcnt})}

\def\qt#1{\qquad\text{#1}}

\def\argmin{\mathop{\rm argmin}}
\def\argmax{\mathop{\rm argmax}}


\setlength{\parskip}{1.4 \medskipamount}
%\sloppy
%\linespread{1.3}

\begin{document}

\title{STAT 238 - Bayesian Statistics  \\
  Lecture Thirty Three} 
\subtitle{Spring 2026, UC Berkeley}

\author{Aditya Guntuboyina}


\date{20 April 2026}

\maketitle

\section{Hamiltonian Monte Carlo}

In HMC, the proposal $y$ is generated by solving the following ODE:
\begin{align}\label{hmc_ode}
  x(0) = x ~~~~ \dot{x}(0) = z ~~~~ \ddot{x}(t) = \nabla \log
  \pi(x(t)) ~~~~ y = x(\sigma). 
\end{align}
Note that if $\nabla \log \pi(x(t))$ is replaced by $\nabla \log
\pi(x)$, then it is easy to check that $y$ equals the MALA proposal $y
= x + (\sigma^2/2)\nabla \log \pi(x) + \sigma z$.

The ODE \eqref{hmc_ode} can be written in the following alternative
first order form:
\begin{align}\label{hmc_firstorder}
  \begin{pmatrix}
    \dot{x}(t) \\ \dot{v}(t)
  \end{pmatrix} =
  \begin{pmatrix}
    v(t) \\ \nabla \log \pi(x(t)) 
  \end{pmatrix} ~\text{with initialization}~
  \begin{pmatrix}
    x(0) \\ v(0)
  \end{pmatrix} =
  \begin{pmatrix}
    x \\ z
  \end{pmatrix}
\end{align}
It is easy to see that \eqref{hmc_ode} and \eqref{hmc_firstorder} are
equivalent. Here is another common way of writing
\eqref{hmc_firstorder}. Let
\begin{align}\label{hamiltonian}
  H(x, v) = H(x_1, \dots, x_d, v_1, \dots, v_d) := -\log \pi(x) +
  \frac{1}{2}\|v\|^2 
\end{align}
This quantity $H(x, v)$ is called the Hamiltonian associated with the
ODE \eqref{hmc_firstorder}. It is then easy to see that
\begin{align*}
  \frac{\partial H}{\partial x} = \left(\frac{\partial H}{\partial
  x_1}, \dots, \frac{\partial H}{\partial x_d} \right)^T = -\nabla
  \log \pi(x)
\end{align*}
and
\begin{align*}
  \frac{\partial H}{\partial v} = \left(\frac{\partial H}{\partial
  v_1}, \dots, \frac{\partial H}{\partial v_d} \right)^T = v. 
\end{align*}
As a result, \eqref{hmc_firstorder} is equivalent to:
\begin{align}\label{hmc_form}
    \begin{pmatrix}
    \dot{x}(t) \\ \dot{v}(t)
  \end{pmatrix} =
  \begin{pmatrix}
    \frac{\partial H}{\partial v} \\ -\frac{\partial H}{\partial x}
  \end{pmatrix} ~\text{with initialization}~
  \begin{pmatrix}
    x(0) \\ v(0)
  \end{pmatrix} =
  \begin{pmatrix}
    x \\ z
  \end{pmatrix}
\end{align}
One important property of this ODE is that the Hamiltonian $H(x(t),
v(t))$ remains constant in $t$ for all times $t$:
\begin{align}\label{ham_const}
  H(x(t), v(t)) = \text{constant} = H(x(0), v(0)) = H(x, z) = -\log
  \pi(x) + \frac{1}{2} \|z\|^2. 
\end{align}
To see \eqref{ham_const}, just note that
\begin{align*}
  \frac{d}{dt} H(x(t), v(t)) &= \sum_{i=1}^d \left( \frac{\partial H}{\partial
                               x_i} \dot{x}_i(t) + \frac{\partial
                               H}{\partial v_i} \dot{v}_i(t) \right) \\
  &= \sum_{i=1}^d \left(\frac{\partial H}{\partial x_i} \frac{\partial
    H}{\partial v_i} + \frac{\partial H}{\partial v_i} \left(-
    \frac{\partial H}{\partial x_i}\right) \right) = \sum_{i=1}^d \left(\frac{\partial H}{\partial x_i} \frac{\partial
    H}{\partial v_i} - \frac{\partial H}{\partial v_i} 
    \frac{\partial H}{\partial x_i} \right) = 0
\end{align*}

\section{Stationarity of $\pi$ for the HMC}

The transition kernel given by \eqref{hmc_ode} satisfies detailed
balance with respect to $\pi$ for every fixed $\sigma > 0$. A
consequence of this is that $\pi$ is invariant (stationary) for this
Markov Chain (for every $\sigma$). We will sketch a proof of this
invariance (and skip the proof of detailed balance).

Stationarity of $\pi$ means that if $x \sim \pi$, then $y$ given by
\eqref{hmc_ode} is also distributed as $\pi$. In order to verify this,
we shall use the following result on density evolutions of ODEs. This
result has been popular in the recent literature on generating
modeling (see e.g., \citet[Equation (5.2.8), Theorem 5.2.2, Section
B.1.2]{lai2025principles}).

\begin{theorem}\label{continuity}
  Consider the ODE 
  \begin{align}\label{ode}
    \dot{S}(x, t) = V(t, S(x, t)) ~~\text{with initial condition $S(x,
    0) = x$}. 
  \end{align}
  Here $\dot{S}(x, t) = \frac{d}{dt}S(x, t)$, $x \in \R^d$ and $S(t,
  \cdot), V(t, \cdot)$ are functions from $\R^d$ to $\R^d$. Suppose
  $\rho_b$ is a density on $\R^d$ and let $\rho(t, \cdot)$ be the 
  density of $S(X, t)$ with $X \sim \rho_b$. Then $\rho(t, x)$
  satisfies the following PDE:
  \begin{align}\label{pdecont}
    \partial_t \rho(t, x) = -\nabla \cdot \left(V(t, x) \rho(t, x)
    \right) ~~ \text{ with $\rho(0, \cdot) = \rho_b$}. 
  \end{align}
  Here $\nabla \cdot$, also known as the divergence, is defined as
  follows. $\nabla \cdot G(x) = \sum_{i=1}^d \frac{\partial}{\partial
    x_i} G_i(x)$ for a function $G(x) = (G_1(x), \dots, G_d(x))$. So
  $\nabla \cdot (V \rho) = \mathrm{div} (V \rho) = \sum_{i=1}^d
  \frac{\partial}{\partial x_i} (V_i(x, t) \rho(x, t))$. 
\end{theorem}

\begin{proof}[Proof of Theorem \ref{continuity}]
  $S(X, t)$ can be interpreted as the position at time $t$ of
  a particle starting at $x$ in a velocity vector field $V(t, x)$,
  formally given by the ODE \eqref{ode}. $\rho(t, \cdot)$ is the pdf
  of $S(X, t)$ with $X \sim \rho_b$.  

  We need to prove that $\rho(t,
  x)$ satisfies the PDE \eqref{pdecont}. Suppose $f$ is a smooth
  function with compact support (i.e., $f \in \C_c^{\infty}$). Then
  $\E f(S(X, t))$ can be written in the following two equivalent ways:
  \begin{equation*}
    \E f(S(X, t)) = \int f(x) \rho(t, x) dx = \int f(S(x, t))
    \rho_b(x) dx. 
  \end{equation*}
  Differentiating both sides of this equality with respect to $t$, we
  obtain
  \begin{equation}\label{temp1_cont}
    \int f(x) \frac{\partial}{\partial t} \rho(t, x) dx = \int
    \left<\dot{S}(x, t), \nabla f(S(x, t)) \right> \rho_b(x) dx
  \end{equation}
  The right hand side above can be simplified using $\dot{S}(x, t) =
  V(t, S(x, t))$ (because $S(x, t)$ satisfies \eqref{ode}) to get
  \begin{align*}
    \int \left<\dot{S}(x, t), \nabla f(S(x, t)) \right> \rho_b(x) dx &=
    \int \left<V(t, S(x, t)), \nabla f(S(x, t)) \right> \rho_b(x) dx \\ &=
                                                                          \int
                                                                          \left<V(t,
n                                                                          x),
                                                                          \nabla
                                                                          f(x)
                                                                          \right>
                                                                          \rho(t,
                                                                          x)
                                                                          dx
    \\ 
    &= \sum_{i=1}^d \int V_i(t, x) \frac{\partial f}{\partial x_i}(x)
      \rho(t, x) dx \\
    &= \sum_{i=1}^d \int V_i(t, x) \rho(t, x) \frac{\partial f}{\partial x_i}(x)dx.
  \end{align*}
Integration by parts (note that we assumed that $f$ has compact
support) and combining with \eqref{temp1_cont} allows us to derive
\begin{align*}
  \int f(x) \frac{\partial}{\partial t} \rho(t, x) dx &= - \sum_{i=1}^d
  \int f(x) \frac{\partial}{\partial x_i} \left(V_i(t, x) \rho(t, x)
                                                        \right) \\
  &= -\int f(x) \sum_{i=1}^d \frac{\partial}{\partial x_i}
    \left(V_i(t, x) \rho(t, x) \right) dx = -\int f(x) \text{div}_x
    \left(V(t, x) \rho(t, x) \right) dx
\end{align*}
Because this holds for every $f$ (with compact support), we deduce the
pde \eqref{pdecont}.
\end{proof}
In the next lecture, we shall see how this result implies that $\pi$
is invariant for the HMC chain. We will basically apply Theorem
\ref{continuity} with $V$ given by:
\begin{align*}
  V(t, x, v) =
  \begin{pmatrix}
    v \\ \nabla \log \pi(x)
  \end{pmatrix} =
  \begin{pmatrix}
    \frac{\partial H}{\partial v} \\ -\frac{\partial H}{\partial x}
  \end{pmatrix}. 
\end{align*}
We shall also discuss the leapfrog discretization of the Hamiltonian
ODE \eqref{hmc_firstorder}, and its Metropolization. 




\begin{thebibliography}{99}

  \bibitem[Lai et~al.(2025)Lai, Song, Kim, Mitsufuji, and Ermon]{lai2025principles}
Lai, C.-H., Song, Y., Kim, D., Mitsufuji, Y., and Ermon, S. (2025).
\newblock The principles of diffusion models.
\newblock \emph{arXiv preprint arXiv:2510.21890}.

%\bibitem[Roberts and Rosenthal(2004)]{RobertsRosenthal2004}
%Roberts, G. O. and Rosenthal, J. S. (2004).
%General state space Markov chains and MCMC algorithms.
%\textit{Probability Surveys}, \textbf{1}, 20--71.

%\bibitem[Meyn and Tweedie(2009)]{MeynTweedie2009}
%Meyn, S. P. and Tweedie, R. L. (2009).
%\textit{Markov Chains and Stochastic Stability}, 2nd ed.
%Cambridge University Press.


  
%\bibitem[Chib and Greenberg(1995)]{chib1995}
%Chib, S. and Greenberg, E. (1995).
%Understanding the Metropolis--Hastings algorithm.
%\textit{The American Statistician}, \textbf{49}, 327--335.


%\bibitem[MacKay and Peto(1995)]{MacKay1995HierarchicalDirichlet}
%D.~J.~C. MacKay and L.~C. Peto,
%\newblock ``A hierarchical Dirichlet language model,''
%\newblock \emph{Natural Language Engineering}, vol.~1,
%no.~3, pp.~289--308, 1995.

%\bibitem[Rubin(1981)]{Rubin1981BayesianBootstrap}
%D.~B. Rubin,
%\newblock ``The Bayesian bootstrap,''
%\newblock \emph{The Annals of Statistics}, vol.~9,
%no.~1, pp.~130--134, 1981.
  
%\bibitem[Fisher(1934)]{Fisher1934}
%R.~A. Fisher,
%\newblock ``Probability, likelihood and quantity of information in the logic of
%uncertain inference,''
%\newblock \emph{Proceedings of the Royal Society of London. Series~A,
%Containing Papers of a Mathematical and Physical Character}, vol.~146,
%no.~856, pp.~1--8, 1934.

%\bibitem[Efron(2010)]{Efron}
%Efron, B. (2010)
%\textit{Large-scale inference: empirical Bayes methods for estimation,
%testing, and prediction}. Cambridge University Press, 2010.

%\bibitem[Shumway and Stoffer(2017)]{ShumwayStoffer}
%Shumway, R. H. and Stoffer, D. S. (2017)
%\textit{Time Series Analysis and Its Applications: With R
%  Examples}. Springer, 4th edition. 
  
%     \bibitem[Jaynes(2003)]{Jaynes} Jaynes, E. T, (2003)
%  \textit{Probability theory: the logic of science}. Cambridge
%  University Press, 2003.
 
%  \bibitem[Sinai(1992)]{Sinai} Sinai, Y. G, (1992)
%  \textit{Probability theory: an introductory course}. Springer
%  Verlag, 1992.  

%\bibitem[{deGroot(1986)}]{DeGrootBlackwell}DeGroot, M. H. (1986)
%       A conversation with David Blackwell.
%\newblock \emph{Statistical Science}, \textbf{1}, 40--53.

%         \bibitem[Mosteller(1987)]{Mosteller} Mosteller, F, (1987)
%  \textit{Fifty challenging problems in probability with
%    solutions}. Courier Corporation, 1987.

%    \bibitem[MacKay(2003)]{MacKay} MacKay, D. J. C., (2003)
%  \textit{Information theory, inference, and learning
%  algorithms}. Cambridge University Press, 2003.

%\bibitem[{Lindley(2013)}]{Lindley}Lindley, D. V. (2013)
%   \textit{Understanding Uncertainty}. John Wiley and sons, 2013.


\end{thebibliography}







\end{document}

%%% Local Variables:
%%% mode: latex
%%% TeX-master: t
%%% End:
